% sec_conclusion.tex — conclusion
\section{Conclusion}\label{sec:conclusion}

We have proposed a substrate-neutral, instantiation-relative measure of
autonomy: the degree to which
behavior is determined by the agent's own normative structure, decomposed into
five channels behind a viability gate, each channel carrying an estimator, a
canonical ratio scale, and a mechanism-labeled zero; the scale is comparable
within a layer at a fixed instantiation, and cross-layer readings are reported
as profiles. The five-channel aggregation rule is
characterized by the stated axioms---forced relative to them, while they are
defended normative commitments rather than measurement-theoretic
necessities---a representation theorem with two derivations
(independent toolchains; the invariance axioms shared), whose minimal
sufficient set is the invariance trio AX4/AX6b/AX8, with zero annihilation
and biseparability following as consequences rather than independent
commitments. The viability gate is an explicit
multiplicative design outside the theorem's hypotheses, defended by its
common-mode semantics rather than absorbed into the derivation. The
theorem suite answers the questions a governance practice actually asks:
which channel sees which intervention (T2, with multi-channel necessity as a
theorem); what manipulation does and what it looks like from inside (T3, with
the hallucination index); where an attacker will strike first and how a
defender allocates (T4, the highest-leverage channel); what total loss means
disjunctively (T5); how much autonomy is retained under continual adversarial
pressure within the registered damage model, and why that margin
comes from repair capacity rather than detection (T6, with the sharp
blind-zone prediction); what delegation costs and when it is free (T7); and
why compelling autonomy is self-undermining within this taxonomy (T8). The
identifiability analysis is
conducted with the same rigor: certified bounds where certification is
possible---one-sided where only one side is certifiable---and a proved
impossibility where no bound exists (implanted endorsement is behaviorally
indistinguishable from endogenous endorsement), making full causal
logging an architectural necessity rather than a convenience. Every claim
carried a falsification window; in $4{,}620$ benchmark runs---an
implementation-level consistency check, not external validation---no
registered window was violated, and the one adverse arm (adaptive
adversaries against the black-box bound) calibrated the fragility of the
bound's premise. These verdicts are internal-consistency results, not
construct validation: construct validity remains unestablished until the
external program below is run (Section~\ref{sec:construct}). The
framework's own falsifier is external: cross-substrate
profiles that fail to predict independent autonomy assessments would
disconfirm the channel choice, not the algebra. The next milestone is stated
as a plan, not a result: the pre-registered independent-substrate port
(protocol outline and power analysis in Section~\ref{sec:threats}(xi))
begins with the MiniGrid-style proof-of-concept---the immediate follow-up
experiment, targeted for the next submission cycle---with the LLM tool-use
harness and the human-side construct-validity study behind it. We claim no indicator
neutrality: the framework's objectivity is procedural---every commitment is
registered, priced, and exposed to a stated falsifier
(Section~\ref{sec:objectivity}).

The framework does not adjudicate whose values autonomy serves, nor whether
machines feel anything; it measures the structures those questions are about.
If autonomy certificates are to mean anything \citep{feng2025}, they must be
issued over causal-layer audit, because no purely behavioral profile can
separate persuasion from hijack. We offer the definitions, the axioms, the
theorems, and the benchmark to make that audit executable.
